\begin{lemma}[The subordination representation of the pseudoinverse $L^{-1/2}\Pi_0$]
Let $f \in L^2(\gamma_n)$, let $\Pi_0 = L^0$ be the mean-zero projection and
$L^{-1/2}\Pi_0$ the pseudoinverse, both as in \noderef{FracPow}, and let $P_s$ be
the Ornstein--Uhlenbeck semigroup of \noderef{OUSemigroup}.  Then:
\begin{enumerate}
\item the $L^2(\gamma_n)$-valued function
$s \mapsto s^{-1/2}\,P_s\Pi_0 f$ is Bochner integrable on $(0,\infty)$; and
\item
\[
  L^{-1/2}\Pi_0 f
    = \frac{1}{\sqrt\pi}\int_0^\infty P_s\Pi_0 f \,\frac{ds}{\sqrt s},
\]
the right-hand side being the Bochner integral of clause (1), taken in
$L^2(\gamma_n)$.
\end{enumerate}

This is the reference paper's displayed convergent Bochner integral for the
pseudoinverse.  The paper writes the integrand as $P_sf - \mathbb{E}(f)\mathbf1$
and remarks that this equals $P_s\Pi_0 f$; the form used here is the latter,
which is the one the spectral description of $\Pi_0$ in \noderef{FracPow} makes
available directly, and the identification of the two forms is neither asserted
nor used below.  The projection is what makes clause (1) true at all: without
removing the constant mode the integrand has $L^2$ norm bounded below by
$|\mathbb{E}(f)|$ for every $s$, and $\int_1^\infty s^{-1/2}ds = \infty$, so the
integral would diverge at infinity whenever $\mathbb{E}(f) \neq 0$ --- exactly
the point the paper makes.

The statement is for every $f \in L^2(\gamma_n)$ with no domain hypothesis, which
is legitimate because $L^{-1/2}\Pi_0$ is everywhere defined: the exponent
$-1/2$ is $\leq 0$, so every element of $L^2(\gamma_n)$ lies in
$\operatorname{Dom}(L^{-1/2})$.
\end{lemma}
\begin{proof}
Write $c_\alpha = \langle f, h_\alpha\rangle$ for the Hermite coefficients of $f$
relative to the basis of \noderef{HermiteBasis}, and recall from that node that
the coefficient map $v \mapsto (v_\alpha)_\alpha$ is an isometric isomorphism of
$L^2(\gamma_n)$ onto $\ell^2(\mathbb{N}_0^n)$; in particular it is injective, and
Parseval's identity $\|v\|_2^2 = \sum_\alpha |v_\alpha|^2$ holds.

\emph{Step 0: the coefficients of $\Pi_0 f$ and of $P_s\Pi_0 f$.}
By \noderef{FracPow} at $s = 0$, the element $\Pi_0 f$ is the element of
$L^2(\gamma_n)$ whose Hermite coefficients are $0$ at $\alpha = 0$ and
$|\alpha|^0 c_\alpha = c_\alpha$ at $\alpha \neq 0$.  By
\noderef{OUSemigroupSpectral}, valid for every $s \geq 0$, the coefficients of
$P_s\Pi_0 f$ are therefore
\begin{equation*}
  (P_s\Pi_0f)_\alpha =
  \begin{cases}
    0, & \alpha = 0, \\
    e^{-s|\alpha|}c_\alpha, & \alpha \neq 0 .
  \end{cases}
  \tag{1}
\end{equation*}

\emph{Step 1: the exponential norm bound.}
Let $s \geq 0$.  By \noderef{OneLeMultiIndexSum} we have $|\alpha| \geq 1$ for
every $\alpha \neq 0$, so $-s|\alpha| \leq -s$ and hence
$e^{-s|\alpha|} \leq e^{-s}$, the exponential being increasing.  By $(1)$ and
Parseval,
\begin{equation*}
  \|P_s\Pi_0f\|_2^2 = \sum_{\alpha\neq0} e^{-2s|\alpha|}c_\alpha^2
    \leq e^{-2s}\sum_{\alpha\neq0}c_\alpha^2
    \leq e^{-2s}\sum_{\alpha}c_\alpha^2 = e^{-2s}\|f\|_2^2 ,
\end{equation*}
all terms being nonnegative.  Taking square roots,
\begin{equation*}
  \|P_s\Pi_0f\|_2 \leq e^{-s}\,\|f\|_2 \qquad (s \geq 0). \tag{2}
\end{equation*}

\emph{Step 2: the scalar integral.}
For every $\lambda > 0$,
\begin{equation*}
  \int_0^\infty s^{-1/2}e^{-\lambda s}\,ds = \sqrt{\pi/\lambda} . \tag{3}
\end{equation*}
Indeed, substitute $s = u^2$, a $C^1$ increasing bijection of $(0,\infty)$ onto
itself with $ds = 2u\,du$ and $s^{-1/2} = u^{-1}$ for $u > 0$; the integrand is
nonnegative and measurable, so the change of variables is valid in $[0,+\infty]$
and gives
\[
  \int_0^\infty s^{-1/2}e^{-\lambda s}\,ds
    = \int_0^\infty u^{-1}e^{-\lambda u^2}\,2u\,du
    = 2\int_0^\infty e^{-\lambda u^2}\,du
    = \int_{\mathbb{R}} e^{-\lambda u^2}\,du
    = \sqrt{\pi/\lambda} ,
\]
the third equality because $u \mapsto e^{-\lambda u^2}$ is even, and the fourth
by the scalar Gaussian integral of the shared context.  In particular $(3)$ with
$\lambda = 1$ gives $\int_0^\infty s^{-1/2}e^{-s}ds = \sqrt\pi < \infty$.

\emph{Step 3: continuity in $s$, hence strong measurability.}
We check that $s \mapsto P_s\Pi_0f$ is continuous from $[0,\infty)$ to
$L^2(\gamma_n)$.  Let $s \geq 0$ and $\varepsilon > 0$.  Since
$\sum_\alpha c_\alpha^2 = \|f\|_2^2 < \infty$, choose a finite set
$S \subseteq \mathbb{N}_0^n$ with $\sum_{\alpha\notin S}c_\alpha^2 < \varepsilon^2/8$.
By $(1)$ and Parseval, for $t \geq 0$,
\[
  \|P_s\Pi_0f - P_t\Pi_0f\|_2^2
    = \sum_{\alpha\neq0}\bigl(e^{-s|\alpha|}-e^{-t|\alpha|}\bigr)^2 c_\alpha^2 .
\]
Split the sum at $S$.  For $\alpha \notin S$ use
$|e^{-s|\alpha|}-e^{-t|\alpha|}| \leq 2$, since both exponentials lie in $[0,1]$;
that part is at most $4\sum_{\alpha\notin S}c_\alpha^2 < \varepsilon^2/2$.  The
remaining sum is finite, and each of its finitely many terms is a continuous
function of $t$ vanishing at $t = s$, so there is $\delta > 0$ with that part
$< \varepsilon^2/2$ whenever $|t-s| < \delta$.  Hence
$\|P_s\Pi_0f - P_t\Pi_0f\|_2 < \varepsilon$ for $|t-s| < \delta$, which is the
claimed continuity.  Consequently $s \mapsto s^{-1/2}P_s\Pi_0f$ is continuous on
$(0,\infty)$, $s \mapsto s^{-1/2}$ being continuous and nonvanishing there; a
continuous function into the separable Banach space $L^2(\gamma_n)$ --- separable
because \noderef{HermiteBasis} exhibits a countable orthonormal basis --- is
strongly measurable.

\emph{Step 4: clause (1).}
By Step 3 the integrand is strongly measurable on $(0,\infty)$, and by $(2)$ its
norm satisfies
$\|s^{-1/2}P_s\Pi_0f\|_2 \leq s^{-1/2}e^{-s}\|f\|_2$, whose integral over
$(0,\infty)$ is $\sqrt\pi\,\|f\|_2 < \infty$ by Step 2.  Bochner integrability
of a function into a Banach space \emph{means} strong measurability together
with an integrable norm --- this is the definition of the notion, not a
criterion imported from elsewhere --- so the function is Bochner integrable on
$(0,\infty)$.  This is clause (1), and
it makes the integral in clause (2) a well-defined element of $L^2(\gamma_n)$.

\emph{Step 5: the coefficients of the integral.}
Fix $\alpha \in \mathbb{N}_0^n$ and let $\Lambda_\alpha(v) = \langle v,h_\alpha\rangle$.
By Cauchy--Schwarz (shared context) $\Lambda_\alpha$ is a continuous linear
functional on $L^2(\gamma_n)$: it is linear in $v$, and
$|\Lambda_\alpha(v)| \leq \|v\|_2\|h_\alpha\|_2 = \|v\|_2$ by orthonormality of
the basis (\noderef{HermiteBasis}).  Writing
$I = \int_0^\infty s^{-1/2}P_s\Pi_0f\,ds$, apply mathlib's
\texttt{ContinuousLinearMap.integral\_comp\_commSL}, which commutes a continuous
linear map with a Bochner integral, with $X = (0,\infty)$ carrying Lebesgue
measure, with $E = L^2(\gamma_n)$ --- a real Banach space, being a Hilbert
space --- with the integrand $F(s) = s^{-1/2}P_s\Pi_0f$, Bochner integrable on
$(0,\infty)$ by Step 4, and with the functional $\Lambda_\alpha$.  Hence
\[
  \Lambda_\alpha(I) = \int_0^\infty s^{-1/2}\,(P_s\Pi_0f)_\alpha \, ds .
\]
If $\alpha = 0$ the integrand vanishes identically by $(1)$, so
$\Lambda_0(I) = 0$.  If $\alpha \neq 0$ then $|\alpha| \geq 1 > 0$ by
\noderef{OneLeMultiIndexSum}, and $(1)$ together with $(3)$ at
$\lambda = |\alpha|$ gives
\[
  \Lambda_\alpha(I) = c_\alpha\int_0^\infty s^{-1/2}e^{-s|\alpha|}\,ds
    = c_\alpha\sqrt{\pi/|\alpha|} = \sqrt\pi\,|\alpha|^{-1/2}c_\alpha ,
\]
the scalar integral being absolutely convergent by Step 2, which is what lets it
be pulled out of the pairing.  Multiplying by $(\sqrt\pi)^{-1}$, the element
$(\sqrt\pi)^{-1}I$ has Hermite coefficients $0$ at $\alpha = 0$ and
$|\alpha|^{-1/2}c_\alpha$ at $\alpha \neq 0$.

\emph{Step 6: conclusion.}
Since $-1/2 \leq 0$, clause (1) of \noderef{FracPowLinear} at $s = -1/2$ gives
$f \in \operatorname{Dom}(L^{-1/2})$, so by \noderef{FracPow} the element
$L^{-1/2}\Pi_0 f$ is the element of $L^2(\gamma_n)$ whose Hermite coefficients
are $0$ at $\alpha = 0$ and $|\alpha|^{-1/2}c_\alpha$ at $\alpha \neq 0$.  By
Step 5 the element $(\sqrt\pi)^{-1}I$ has the same coefficients, and the
coefficient map is injective by \noderef{HermiteBasis}.  Therefore
$L^{-1/2}\Pi_0 f = (\sqrt\pi)^{-1}I$, which is clause (2).
\end{proof}
